\documentclass[10pt,a4paper]{amsart}
\usepackage{amsmath,amssymb,amsthm,graphicx}
% Retained preamble for the Stacks Project chapter snapshot in this directory.
% The chapter's own preamble.tex is a copy of the shared Stacks Project preamble; the block
% below is sliced MECHANICALLY, line for line, from it. Nothing was added, removed or
% redefined; the wrapper class, the xr-hyper/hyperref wiring and the externaldocument list
% are omitted because the assembler supplies the document class.
%
% lines taken from preamble.tex: 3, 5, 9, 11, 12, 156, 157, 158, ... (51 lines total)
%
% package coverage restored from the chapter's own preamble (2026-10-13): xy, xy-2cell

\usepackage[all]{xy}
\xyoption{2cell}
\UseAllTwocells

\newenvironment{reference}{\comment}{\endcomment}
\newenvironment{slogan}{\comment}{\endcomment}
\newenvironment{history}{\comment}{\endcomment}
\theoremstyle{plain}
\newtheorem{theorem}[subsection]{Theorem}
\newtheorem{proposition}[subsection]{Proposition}
\newtheorem{lemma}[subsection]{Lemma}
\theoremstyle{definition}
\newtheorem{definition}[subsection]{Definition}
\newtheorem{example}[subsection]{Example}
\newtheorem{exercise}[subsection]{Exercise}
\newtheorem{situation}[subsection]{Situation}
\theoremstyle{remark}
\newtheorem{remark}[subsection]{Remark}
\newtheorem{remarks}[subsection]{Remarks}
\numberwithin{equation}{subsection}
\def\lim{\mathop{\mathrm{lim}}\nolimits}
\def\colim{\mathop{\mathrm{colim}}\nolimits}
\def\Spec{\mathop{\mathrm{Spec}}}
\def\Hom{\mathop{\mathrm{Hom}}\nolimits}
\def\Ext{\mathop{\mathrm{Ext}}\nolimits}
\def\SheafHom{\mathop{\mathcal{H}\!\mathit{om}}\nolimits}
\def\SheafExt{\mathop{\mathcal{E}\!\mathit{xt}}\nolimits}
\def\Sch{\mathit{Sch}}
\def\Mor{\mathop{\mathrm{Mor}}\nolimits}
\def\Ob{\mathop{\mathrm{Ob}}\nolimits}
\def\Sh{\mathop{\mathit{Sh}}\nolimits}
\def\NL{\mathop{N\!L}\nolimits}
\def\CH{\mathop{\mathrm{CH}}\nolimits}
\def\proetale{{pro\text{-}\acute{e}tale}}
\def\etale{{\acute{e}tale}}
\def\QCoh{\mathit{QCoh}}
\def\Ker{\mathop{\mathrm{Ker}}}
\def\Im{\mathop{\mathrm{Im}}}
\def\Coker{\mathop{\mathrm{Coker}}}
\def\Coim{\mathop{\mathrm{Coim}}}
\DeclareMathSymbol{\boxtimes}{\mathbin}{AMSa}{"02}
\def\QCohstack{\mathcal{QC}\!\mathit{oh}}
\def\Cohstack{\mathcal{C}\!\mathit{oh}}
\def\Spacesstack{\mathcal{S}\!\mathit{paces}}
\def\Quotfunctor{\mathrm{Quot}}
\def\Hilbfunctor{\mathrm{Hilb}}
\def\Curvesstack{\mathcal{C}\!\mathit{urves}}
\def\Polarizedstack{\mathcal{P}\!\mathit{olarized}}
\def\Complexesstack{\mathcal{C}\!\mathit{omplexes}}
\def\Pic{\mathop{\mathrm{Pic}}\nolimits}
\def\Picardstack{\mathcal{P}\!\mathit{ic}}
\def\Picardfunctor{\mathrm{Pic}}
\def\Deformationcategory{\mathcal{D}\!\mathit{ef}}
\usepackage{cleveref}
\crefname{originallemma}{original lemma}{original lemmas}
\Crefname{originallemma}{Original lemma}{Original lemmas}
\makeatletter
\newlabel{algebra-lemma-characterize-zero-local}{{00HN}{source}{}{sourcepointer-1}{}}
\newlabel{algebra-lemma-characterize-zero-local@cref}{{[originallemma][0][]00HN}{[1][0][]source}{}{}{}}
\makeatother
\crefname{originallemma}{original lemma}{original lemmas}
\Crefname{originallemma}{Original lemma}{Original lemmas}
\makeatletter
\newlabel{algebra-lemma-etale-makes-quasi-finite-finite-one-prime}{{00UJ}{source}{}{sourcepointer-2}{}}
\newlabel{algebra-lemma-etale-makes-quasi-finite-finite-one-prime@cref}{{[originallemma][0][]00UJ}{[1][0][]source}{}{}{}}
\makeatother
\crefname{originallemma}{original lemma}{original lemmas}
\Crefname{originallemma}{Original lemma}{Original lemmas}
\makeatletter
\newlabel{algebra-lemma-local-flat-ff}{{00HR}{source}{}{sourcepointer-3}{}}
\newlabel{algebra-lemma-local-flat-ff@cref}{{[originallemma][0][]00HR}{[1][0][]source}{}{}{}}
\makeatother
\crefname{originallemma}{original lemma}{original lemmas}
\Crefname{originallemma}{Original lemma}{Original lemmas}
\makeatletter
\newlabel{equation-bc-dualizing}{{0BVB}{source}{}{sourcepointer-4}{}}
\newlabel{equation-bc-dualizing@cref}{{[originallemma][0][]0BVB}{[1][0][]source}{}{}{}}
\makeatother
\crefname{originallemma}{original lemma}{original lemmas}
\Crefname{originallemma}{Original lemma}{Original lemmas}
\makeatletter
\newlabel{lemma-bc-map-dualizing}{{0BVC}{source}{}{sourcepointer-5}{}}
\newlabel{lemma-bc-map-dualizing@cref}{{[originallemma][0][]0BVC}{[1][0][]source}{}{}{}}
\makeatother
\crefname{originallemma}{original lemma}{original lemmas}
\Crefname{originallemma}{Original lemma}{Original lemmas}
\makeatletter
\newlabel{lemma-dualizing-flat-base-change}{{0BT3}{source}{}{sourcepointer-6}{}}
\newlabel{lemma-dualizing-flat-base-change@cref}{{[originallemma][0][]0BT3}{[1][0][]source}{}{}{}}
\makeatother
\crefname{originallemma}{original lemma}{original lemmas}
\Crefname{originallemma}{Original lemma}{Original lemmas}
\makeatletter
\newlabel{lemma-dualizing-product}{{0BT5}{source}{}{sourcepointer-7}{}}
\newlabel{lemma-dualizing-product@cref}{{[originallemma][0][]0BT5}{[1][0][]source}{}{}{}}
\makeatother
\title{Stacks, Discriminants: lemma-dualizing-base-change-of-flat}
\author{The Stacks Project Authors}
\date{stacks-project:discriminant.tex:local-snapshot-20261012; GNU FDL 1.2 or later}
\begin{document}
\maketitle
\noindent\textit{Editorial note.} Statement span is lines 589-593 of sources/stacks-discriminant/pinned-original-source.tex; proof span is lines 595-646, the complete author proof (50 body lines).\par\medskip
\noindent\textit{Editorial note.} Navigation stubs: 7 cited-but-not-included labels are supplied as non-load-bearing pointers with their pinned tags where the snapshot tag index has them.\par\medskip
\noindent\textit{Editorial note.} extra\_packages is empty on purpose: the retained preamble already loads xy with options and a second declaration clashes in the assembler wrapper.\par\medskip
\noindent\textit{Editorial note.} Text is retained verbatim, including the author's own typos and placeholders.\par\medskip
\noindent\textit{Editorial note.} Queue index 6 of queue\_discriminant.json (pb=2323, stmt 589-593, proof 595-646).\par\medskip
\noindent\textit{Allowed standard prerequisites (editorial).}\begin{itemize}
\item Explicit prerequisite boundary. The cited inputs are prerequisites; the argument in the retained span is not replaced by any cited result.
\end{itemize}\medskip
\ifdefined\hypertarget\hypertarget{sourcepointer-1}{}\fi
\noindent\textit{Source reference.} Original lemma 00HN: sources/stacks-discriminant/pinned-original-source.tex, source lines 646--646. Cited by the retained statement or proof but not part of the collected spans; supplied only as a navigation pointer. The argument does not depend on its content.\par\medskip
\ifdefined\hypertarget\hypertarget{sourcepointer-2}{}\fi
\noindent\textit{Source reference.} Original lemma 00UJ: sources/stacks-discriminant/pinned-original-source.tex, source lines 646--646. Cited by the retained statement or proof but not part of the collected spans; supplied only as a navigation pointer. The argument does not depend on its content.\par\medskip
\ifdefined\hypertarget\hypertarget{sourcepointer-3}{}\fi
\noindent\textit{Source reference.} Original lemma 00HR: sources/stacks-discriminant/pinned-original-source.tex, source lines 646--646. Cited by the retained statement or proof but not part of the collected spans; supplied only as a navigation pointer. The argument does not depend on its content.\par\medskip
\ifdefined\hypertarget\hypertarget{sourcepointer-4}{}\fi
\noindent\textit{Source reference.} Original lemma 0BVB: sources/stacks-discriminant/pinned-original-source.tex, source lines 414--416. Cited by the retained statement or proof but not part of the collected spans; supplied only as a navigation pointer. The argument does not depend on its content.\par\medskip
\ifdefined\hypertarget\hypertarget{sourcepointer-5}{}\fi
\noindent\textit{Source reference.} Original lemma 0BVC: sources/stacks-discriminant/pinned-original-source.tex, source lines 435--441. Cited by the retained statement or proof but not part of the collected spans; supplied only as a navigation pointer. The argument does not depend on its content.\par\medskip
\ifdefined\hypertarget\hypertarget{sourcepointer-6}{}\fi
\noindent\textit{Source reference.} Original lemma 0BT3: sources/stacks-discriminant/pinned-original-source.tex, source lines 450--454. Cited by the retained statement or proof but not part of the collected spans; supplied only as a navigation pointer. The argument does not depend on its content.\par\medskip
\ifdefined\hypertarget\hypertarget{sourcepointer-7}{}\fi
\noindent\textit{Source reference.} Original lemma 0BT5: sources/stacks-discriminant/pinned-original-source.tex, source lines 493--497. Cited by the retained statement or proof but not part of the collected spans; supplied only as a navigation pointer. The argument does not depend on its content.\par\medskip
\begin{lemma}
\label{lemma-dualizing-base-change-of-flat}
If $A \to B$ is flat, then the base change map (\ref{equation-bc-dualizing})
induces an isomorphism $\omega_{B/A} \otimes_B B_1 \to \omega_{B_1/A_1}$.
\end{lemma}

\begin{proof}
If $A \to B$ is finite flat, then $B$ is finite locally free as an $A$-module.
In this case $\omega_{B/A} = \Hom_A(B, A)$ is the dual finite
locally free $A$-module and formation of this module commutes
with arbitrary base change which proves the lemma in this case.
In the next paragraph we reduce the general (quasi-finite flat)
case to the finite flat case just discussed.

\medskip\noindent
Let $\mathfrak q_1 \subset B_1$ be a prime. We will show that the
localization of the map at the prime $\mathfrak q_1$ is an isomorphism, which
suffices by Algebra, Lemma \ref{algebra-lemma-characterize-zero-local}.
Let $\mathfrak q \subset B$ and $\mathfrak p \subset A$ be the prime
ideals lying under $\mathfrak q_1$. By
Algebra, Lemma \ref{algebra-lemma-etale-makes-quasi-finite-finite-one-prime}
we can find an \'etale ring map $A \to A'$ and a prime
ideal $\mathfrak p' \subset A'$ lying over $\mathfrak p$
such that $\kappa(\mathfrak p') = \kappa(\mathfrak p)$ and
such that
$$
B' = B \otimes_A A' = C \times D
$$
with $A' \to C$ finite and such that the unique prime $\mathfrak q'$
of $B \otimes_A A'$ lying over $\mathfrak q$ and $\mathfrak p'$
corresponds to a prime of $C$. Set $A'_1 = A' \otimes_A A_1$ and
consider the base change maps
(\ref{equation-bc-dualizing}) for the ring maps
$A \to A' \to A'_1$ and $A \to A_1 \to A'_1$ as in the diagram
$$
\xymatrix{
\omega_{B'/A'} \otimes_{B'} B'_1 \ar[r] & \omega_{B'_1/A'_1} \\
\omega_{B/A} \otimes_B B'_1 \ar[r] \ar[u] &
\omega_{B_1/A_1} \otimes_{B_1} B'_1 \ar[u]
}
$$
where $B' = B \otimes_A A'$, $B_1 = B \otimes_A A_1$, and
$B_1' = B \otimes_A (A' \otimes_A A_1)$. By
Lemma \ref{lemma-bc-map-dualizing} the diagram commutes. By
Lemma \ref{lemma-dualizing-flat-base-change}
the vertical arrows are isomorphisms.
As $B_1 \to B'_1$ is \'etale and hence flat it suffices
to prove the top horizontal arrow is an isomorphism after localizing
at a prime $\mathfrak q'_1$ of $B'_1$ lying over $\mathfrak q$
(there is such a prime and use
Algebra, Lemma \ref{algebra-lemma-local-flat-ff}).
Thus we may assume that $B = C \times D$ with $A \to C$
finite and $\mathfrak q$ corresponding to a prime of $C$.
In this case the dualizing module $\omega_{B/A}$ decomposes
in a similar fashion (Lemma \ref{lemma-dualizing-product})
which reduces the question
to the finite flat case $A \to C$ handled above.
\end{proof}

\end{document}
